
\section{別名}
       耶琳- 就是基撒崙\\
       耶布斯- 就是耶路撒冷\\
       底璧- 從前名叫基列西弗\\
       基列亞巴-  就是希伯崙\\
       加略希斯崙- 就是夏瑣\\
       基列薩拿- 就是底璧\\
       基列巴力- 就是基列耶琳\\
       路斯- 就是伯特利\\
       利善- 改名為但\\
       別是巴- 或名示巴\\


\section{地名}
\subsection{A}
\begin{itemize} 
\itemsep-.25em 
\item ABILENE Ἀβιληνή (Region) \\
Probably from Hebrew אָבֵל ('avel) meaning "meadow, grassy place". This is the name of a place briefly mentioned in the New Testament.
\item Achor valley 亚割谷  עָכ֔וֹר  亚割就是连累trouble的意思（約7:24-26）  עֲכַרְתָּ֔נוּ: have you troubled. יַעְכֳּרְךָ֥:shall trouble. And Joshua said, Why have you troubled ( עֲכַרְתָּ֔נוּ) us? the LORD shall trouble you (יַעְכֳּרְךָ֥ ) this day\\
\item APHRA רפע f \\
Meaning uncertain; possibly a variant of AFRA (1), or possibly a variant of Aphrah, a biblical place name meaning "dust". 
\item ARARAT אֲרָרָט (Mountain) \\
From the name of the ancient kingdom of Urartu. This is the name of a mountain in Turkey (formerly part of Armenia), the place where Noah's Ark came to rest according to the Old Testament.
\item ARIEL אֲרִיאֵל m and f \\
Means "lion of God" in Hebrew, from אֲרִי ('ari) meaning "lion" and אֵל ('el) meaning "God". In the Old Testament it is used as another name for the city of Jerusalem. 
\begin{itemize} 
\itemsep-.25em 
\item The name given to Jerusalem in Isaiah 29:1–2, 7, where God will bring distress upon Ariel, and will make her like an ariel (for meaning, see below). Ariel in this sense is probably connected with the form Erellam ( אֶרְאֵל perhaps a hero) in Isaiah 33:7, understood as the plural form arielim ("Jerusalemites"), parallel to "messengers of Shalom" (i.e., of Jerusalem; cf. Gen. 14:18; Ps. 76:3).
\item A cultic object in Ezekiel 43:15–16, where it occurs in the forms ariel羚羊之类 and harel. This is apparently an altar祭坛;供坛 hearth炉床 superimposed重叠 upon the base of the altar, having horns at its four corners. Alternatively, it may be viewed as the top two sections of a three-tiered altar, again with the function of a hearth. This usage has been connected by some with the ʾrʾl dwdh which Mesha of Moab dragged before Chemosh from a captured town (Mesha Stele, 1:12). It may also be connected with (II Samuel 23:20 = I Chronicles 11:22) "The two Ariels of Moab." 
\item One of the chief men summoned by Ezra in Ezra 8:16 (cf. also Gen. 46:16; Num. 26:17). The etymology语源 of this word is the subject of some dispute. Three principal modes of interpretation have been proposed: 
\begin{itemize} 
\itemsep-.25em 
\item from ari-el, "lion of God" or "Great Lion." This is the most probable derivation for the personal name in Ezra; 
\item from a posited root ari, "to burn," with lamed希伯莱语第十二个字母 afformative确定, thus meaning "hearth," similar to the Arabic ʿiratun, "hearth"
\item as a loanword from the Akkadian arallû-, the name for the netherworld and allegedly the world mountain. (In this view the altar is understood as a miniature ziggurat, which is taken to be the symbol of the world mountain.) However, arallû does not mean "mountain." In addition, the Akkadian, a loanword from Sumerian, would have not shown up in Hebrew in the form attested. Regardless of the ultimate derivation of the word, the meaning of Isaiah 29:1–2 seems to be that Jerusalem, here (prophetically?) called Ariel, is to become like the altar, i.e., a scene of holocaust.
\end{itemize} 
\item ARIMATHEA Ἁριμαθαία (Settlement) \\
From Greek Ἁριμαθαία (Harimathaia), of unknown meaning. In the New Testament this is the home town of Joseph of Arimathea. The town has not been positively identified, though רָמָתַיִם (Ramatayim) or רָמָה (Ramah) near Jerusalem has been proposed.
\end{itemize} 
\end{itemize} 

\subsection{B}
\begin{itemize} 
\itemsep-.25em
\item 別迦摩
\item 別加 (Perga)
要塞 Citadel, 土質的,世俗的 Earthy �Murtana,位於今日土耳其的南部,Cestrus 河西岸,距海約 13 公里,距他的港埠亞大利 (今日的 Antalya 大城) 約 17 公里,距安提阿約 155 公里,因有來自彼西底高原的水流灌溉,故土地肥沃,民豐物富,在新約時代,它曾是羅馬帝國旁非利亞省的省城,在現有的遺跡中,尚有一座羅馬式的戲院,估計可容一萬八千餘人,另有華麗的浴池,有頂棚的長街、田徑場等等,都足以表明它曾是一個人口稠密的城市,該城又以奉拜女神 Artemis (Diana) 而聞名。
徒 13:13 保羅等從帕弗開船,來到旁非利亞的別加,然後去彼西底的安提阿,14:24 回程時再到別加講道後去亞大利。  
\item Baalah 巴拉 בַּעֲלָ֔ה - 就是基列耶琳 Kirjathjearim  יְעָרִֽים 
\item BABYLON Βαβυλών (Settlement) \\
Greek form of Akkadian (Babili), which appears to mean "gateway of God", from Akkadian (babu) meaning "gate" and (ilu) meaning "God", though it may in fact derive from a non-Semitic language. This was the name of a major city in ancient Mesopotamia, the capital of the Babylonian Empire. It was located in present-day Iraq.
\item BETHANY (Settlement) \\
From Greek Βηθανία (Bethania), which is of uncertain meaning. The first part of the name is derived from Hebrew/Aramaic בַּיִת (bayit) meaning "house". Suggestions for the second part of the name include עָנָה ('anah) leading to "house of affliction" or תְּאֵנָה (te'enah) leading to "house of figs". In the New Testament the town of Bethany is the home of Lazarus and his sisters Mary and Martha.
\item BETHEL בֵּית־אֵל (Settlement) \\
Means "house of God" in Hebrew. In the Old Testament this is a town north of Jerusalem, where Jacob saw his vision of the stairway.
\item BETHLEHEM בֵּית־לֶחֶם (Settlement) \\
Means "house of bread" in Hebrew, from the roots בַּיִת (bayit) meaning "house" and לֶחֶם (lechem) meaning "bread". This is the name of a city in Palestine. It appears in the both the Old Testament and the New Testament, notably as the town where Jesus is born.
\item BEULAH בְּעוּלָ֑ה bə'ūlāh f\\
Means "married" in Hebrew. The name is used in the Old Testament to refer to the land of Israel (Isaiah 62:4). 
\item Bochim  בֹּכִ֑ים  波金：就是哭的意思（士師記2：5）\\
\end{itemize} 
%\end{itemize} 

\subsection{C}
\begin{itemize} 
\itemsep-.25em 
\item CARMEL (Mountain) כַּרְמִיאֵל \\
means "fresh" (planted), or "vineyard" (planted). The word is also commonly used in reference to the Carmelites and their houses or foundations. 
\begin{itemize} 
\itemsep-.25em 
\item Mount Carmel, coastal mountain range in Israel overlooking the Mediterranean Sea
\item Carmel (biblical settlement), an ancient Israelite town in Judea
\end{itemize}
\item CANAAN כְּנָעַן‎  m \\
Meaning unknown. In the Old Testament this is the name of a son of Ham. He is said to be the ancestor of the Canaanite people.
\end{itemize} 

 
\subsection{E}
\begin{itemize} 
\itemsep-.25em 
\item EBENEZER 以便以謝 אֶבֶן הָעֵזֶר  \\
Means "stone of help" in Hebrew. This was the name of a monument erected by Samuel in the Old Testament. 1 Samuel 4:1, 5:1, 7:12. a place of unc. location, also a commemorative纪念的 stone. 撒母耳將一塊石頭立在米斯巴和善的中間，給石頭起名叫以便以謝，說：「到如今耶和華都幫助我們。」(撒上7:12）
\item EDEN עֵדֶן (Region) \\
Possibly from Hebrew עֵדֶן ('eden) meaning "pleasure, delight", or perhaps derived from Sumerian (edin) meaning "plain". According to the Old Testament the Garden of Eden was the place where the first people, Adam and Eve, lived before they were expelled.
\item EDOM אֱדֹם m Genesis 25:30 \\
From Hebrew אָדֹם ('adom) meaning "red". According to the Old Testament, Esau, who is described as having red skin, was given this name because he traded his birthright for a helping of red broth. The bible goes on to tell that Esau was the founder of the ancient nation of Edom, located to the south of the kingdom of Judah.
\item ERECH (Arkevaye אַרְכְּוַי  ( אֶרֶכְ (Settlement) \\
Form of URUK used in the English Old Testament. Ezra 4:9
\end{itemize} 


\subsection{F}
非拉鐵非

\subsection{G}
\begin{itemize} 
\itemsep-.25em 
\item GETHSEMANE (Region) \\
From Γεθσημανί (Gethsemani), the Greek form of an Aramaic place name meaning "oil press". In the New Testament this is the name of the garden where Jesus was arrested, located on the Mount of Olives near Jerusalem.
\item GILEAD גִּלְעָד m \\
From an Old Testament place name meaning "heap of witness" in Hebrew. This was a mountainous region east of the Jordan River. Besides being a place name, it is also borne by people in the Bible.
\item Gilgal 吉甲 גִּלְגָּ֔ל : 就是滾的意思（約5：9）\\
\end{itemize} 

\subsection{H-I}
\begin{itemize} 
\itemsep-.25em 
\item HAVILAH חֲוִילָה m Genesis 2:11,10:7,29,25:18\\
Probably means "to dance, to circle, to twist" in Hebrew. In the Old Testament this is both a place name and a masculine personal name.
\item ISRAEL יִשְׂרָאֵל, Ἰσραήλ \\
From the name of the Old Testament hero Jacob, who was also called ISRAEL, "God strives". 
\item IVAH עִוָּה f 2 Kings 17:24,18:34\\
From the name of a district of Babylon, a city conquered by Assyr. mentioned in the Old Testament.
\end{itemize} 

\subsection{J}
\begin{itemize} 
\itemsep-.25em 
\item Jehovah-Shalom 耶和華沙龍  שלום יהוה  :於是基甸在那裡為耶和華築了一座壇，起名叫耶和華沙龍。這壇在亞比以謝族的俄弗拉，直到如今。(士 6:24）\\ 
\item JERICHO יְרִיחוֹ (Settlement) \\
Meaning uncertain, possibly related to the Hebrew word יָרֵחַ (yareach) meaning "moon", or otherwise to the Hebrew word רֵיחַ (reyach) meaning "fragrant". This is the name of a city in Palestine, mentioned several times in the Old Testament.
\item JERUSALEM (Settlement) \\
From Hebrew יְרוּשָׁלַיִם (Yerushalayim), from an earlier Canaanite form like Urushalim, probably meaning "established by (the god) SHALIM". This is the name of a city in Israel and Palestine. Originally a Canaanite city, it was conquered by the Israelites under King David. 
\item JORDAN יַרְדֵן \\
River that flows between the countries of Jordan and Israel. The river's name in Hebrew is יַרְדֵן (Yarden), and it is derived from יָרַד (yarad) meaning "descend" or "flow down". 
\item JUDAH m \\
From the Hebrew name יְהוּדָה (Yehudah), probably derived from יָדָה (yadah) meaning "praise". In the Old Testament Judah is the fourth of the twelve sons of Jacob by Leah, and the ancestor of the tribe of Judah. An explanation for his name is given in Genesis 29:35. His tribe eventually formed the Kingdom of Judah in the south of Israel. King David and Jesus were among the descendants of him and his wife Tamar.
\end{itemize} 

\subsection{L-P}
\begin{itemize} 
\itemsep-.25em
\item 老底嘉  
\item LEHI לֶחִי (Region) \\
Means "jawbone" in Hebrew. In the Old Testament this is the name of the site where the hero Samson killed 1,000 men using only a donkey's jawbone.
\item Mizpah  מִצְפָּה 米斯巴 撒上 7:12 \\
Hebrew for "watchtower". As mentioned in the biblical story of Jacob and Laban, making a pile of stones marked an agreement between two people, with God as their watching witness.
\item MOAB מוֹאָב m \\
Means "of his father" in Hebrew. In the Old Testament this is the name of a son of Lot. He was the ancestor of the Moabites, a people who lived in the region called Moab to the east of Israel.
\item MORIAH מֹרִיָה (Mountain) Genesis 22:2 \\
Possibly means "seen by YAHWEH" in Hebrew. In the Old Testament this is both the place where Abraham is to sacrifice Isaac and the mountain upon which Solomon builds the temple.
\item NAZARETH נָצְרַת (Settlement) \\
Possibly means "branch" in Hebrew. This is the name of a town in Galilee in Israel. It was the home town of Jesus.
\item NINEVEH (Settlement) \\
From Akkadian (Ninua), possibly related to (nūnu) meaning "fish". When written, the name is usually prefixed with indicating a city. This may have referred to an aspect of Ishtar, as from an early time the city was a center of worship of the goddess. The cuneiform symbols used to represent the city's name depict a fish within a house
\item OPHRAH עָפְרָה m \\
Means "fawn" in Hebrew. In the Old Testament this is the name of both a man mentioned in genealogies and a city in Manasseh
\item PHILADELPHIA Φιλαδέλφεια (Settlement) \\
Means "brotherly love" from Greek φιλέω (phileo) meaning "to love" and ἀδελφός (adelphos) meaning "brother". This was the name of a city in Asia Minor mentioned in Revelation in the New Testament. It is now known as Alaşehir (in Turkey). 
\end{itemize} 

\subsection{S-T}
\begin{itemize} 
\itemsep-.25em
\item 撒狄  
\item SELA סֶלַע (Settlement) 2 Kings 14:7,Isaiah 16:1,42:11 \\
Means "rock" in Hebrew. In the Old Testament this is the name of a city, the capital of Edom. 
\item SHARON שָׁרוֹן (Region) Joshua 12:18,1 Chronicles 5:16\\
Means "plain" in Hebrew, referring to the fertile plain near the coast of Israel.
\item SHEBA m Genesis 10:7,10:28\\
  Means "oath" in Hebrew. This is the name of several characters in the Old Testament. Also in the Bible, this referring to region Ethiopia衣索比亚. The queen of Sheba visited Solomon after hearing of his wisdom.
\item 士每拿  
\item SHILOH שִׁילֹה  Joshua 18:1,8-10\\
From an Old Testament place name possibly meaning "tranquil安静的" in Hebrew. It is also used prophetically in the Old Testament to refer to a person, often understood to be the Messiah (see Genesis 49:10). This may in fact be a mistranslation. 
\item TARAH תָּרַח (Settlement) Genesis 11:24-27\\
A place name (an encampment扎营) used in some versions of the Old Testament. It is identical to the personal name TERAH.
\item TEMAN תֵּימָן (Settlement) Genesis 36:11,15,42;Jeremiah 49:7,20;Ezekiel 25:13\\
  Means "right hand" or "south" in Hebrew. This is an Edomite town in the Old Testament, supposedly named for a grandson of Esau.
\item 推雅推喇
\item 以弗所
\end{itemize}
